\section{AIME I 2000}
        \begin{enumerate}
        	\item (\textit{AIME I 2000}) Find the least positive integer $n$ such that no matter how $10^{n}$ is expressed as the product of any two positive integers, at least one of these two integers contains the digit $0$.


 
 \textbf{Answer: }008



\textbf{Solution}

Let $a$ and $b$ be positive integers such that $ab = 10^n$.

 If neither $a$ nor $b$ contains $0$, then neither $a$ nor $b$ is divisible by $10$ (and thus neither $a$ nor $b$ contains both $2$ and $5$ in its prime factorization).  This is only possible if $\{a, b\} = \{2^n, 5^n\}$.  

 Thus the answer is the smallest integer $n$ such that either $2^n$ or $5^n$ contains $0$.  By inspection, the smallest power of $2$ containing $0$ is $2^{10} = 1024$, while the smallest power of $5$ containing $0$ is $5^8 = 390625$.  Thus $n = \boxed{008}$ is the smallest qualifying value of $n$.

 




	\item (\textit{AIME I 2000}) Let $u$ and $v$ be integers satisfying $0 < v < u$. Let $A = (u,v)$, let $B$ be the reflection of $A$ across the line $y = x$, let $C$ be the reflection of $B$ across the y-axis, let $D$ be the reflection of $C$ across the x-axis, and let $E$ be the reflection of $D$ across the y-axis. The area of pentagon $ABCDE$ is $451$. Find $u + v$.


 
 \textbf{Answer: }021



\textbf{Solution 1}


                    \begin{figure}[H]
                        \centering
                        \includegraphics[width=0.5\textwidth]{Solution/AIME_2000_I_2_Solution_Figure_1.png}
                    \end{figure}
                    Since $A = (u,v)$, we can find the coordinates of the other points: $B = (v,u)$, $C = (-v,u)$, $D = (-v,-u)$, $E = (v,-u)$. If we graph those points, we notice that since the latter four points are all reflected across the x/y-axis, they form a rectangle, and $ABE$ is a triangle. The area of $BCDE$ is $(2u)(2v) = 4uv$ and the area of $ABE$ is $\frac{1}{2}(2u)(u-v) = u^2 - uv$. Adding these together, we get $u^2 + 3uv = u(u+3v) = 451 = 11 \cdot 41$. Since $u,v$ are positive, $u+3v>u$, and by matching factors we get either $(u,v) = (1,150)$ or $(11,10)$. Since $v < u$ the latter case is the answer, and $u+v = \boxed{021}$.

 





\textbf{Solution 2}

We find the coordinates like in the solution above: $A = (u,v)$, $B = (v,u)$, $C = (-v,u)$, $D = (-v,-u)$, $E = (v,-u)$. Then we notice pentagon $ABCDE$ fits into a rectangle of side lengths $(u+v)$ and $(2u)$, giving us two triangles, each with hypotenuse $AB$ and $BE$. First, we can solve for the first triangle. Using the coordinates of $A$ and $B$, we discover the side lengths are both $(u-v)$, so the area of the triangle of hypotenuse $AB$ is $\frac{1}{2}(u-v)^2$. Next, we can solve for the second triangle. Using the coordinates of $A$ and $E$, we discover the side lengths are $(u-v)$ and $(u+v)$, so the area of the triangle of hypotenuse $AE$ is $\frac{1}{2}(u-v)(u+v) = \frac{1}{2}(u^2-v^2)$. Now, let’s subtract the area of these 2 triangles from the rectangle giving us $(u+v)(2u)-\frac{1}{2}(u-v)^2-\frac{1}{2}(u^2-v^2)=451 —> u^2+3uv=451 -> u(u+3v)=451$. Next, we take note of the fact that $u$ and $u+3v$ are both factors of 451, and since both $u$ and $v$ are positive integers, $u+3v$ must be greater than $u$, thus giving us two cases, where either $u=1$ or $u=11$. After trying both, the only working pair of $(u,v)$ where both $u$ and $v$ are integers are $u=11$ and $v=10$, thus meaning $u + v =$ $\boxed{021}$

 





\textbf{Solution 3}

We find the coordinates like in the solution above: $A = (u,v)$, $B = (v,u)$, $C = (-v,u)$, $D = (-v,-u)$, $E = (v,-u)$. Then we apply the \href{https://artofproblemsolving.com/wiki/index.php?title=Shoelace\_Theorem}{Shoelace Theorem}. 
 \[A =  \frac{1}{2}[(u^2 + vu + vu + vu + v^2) - (v^2 - uv - uv - uv -u^2)] = 451\] 
 \[\frac{1}{2}(2u^2 + 6uv) = 451\] 
 \[u(u + 3v) = 451\]

 This means that $(u,v) = (11, 10)$ or $(1,150)$, but since $v < u$, then the answer is $\boxed{021}$

 




	\item (\textit{AIME I 2000}) In the expansion of $(ax + b)^{2000},$ where $a$ and $b$ are relatively prime positive integers, the coefficients of $x^{2}$ and $x^{3}$ are equal. Find $a + b$.


 
 \textbf{Answer: }667



\textbf{Solution}

Using the \href{https://artofproblemsolving.com/wiki/index.php?title=Binomial\_theorem}{binomial theorem}, $\binom{2000}{1998} b^{1998}a^2 = \binom{2000}{1997}b^{1997}a^3 \Longrightarrow b=666a$.

 Since $a$ and $b$ are positive relatively prime integers, $a=1$ and $b=666$, and $a+b=\boxed{667}$.

 Note: The ratio of $\binom{2000}{1997} b^{1997} a^3$ to $\binom{2000}{1998} b^{1998} a^2$ can be calculated by writing them as 
 \[\frac{\frac{2000!}{1997! \cdot 3!}}{\frac{2000!}{1998! \cdot 2!}}.\]Simplifying, we get 
 \[\frac{1998! \cdot 2!}{1997! \cdot 3!} = \frac{1998}{3} = 666.\]

 $666$ is the value for $b$. Since $a=1$, $a+b=\boxed{667}$.

 




	\item (\textit{AIME I 2000}) The diagram shows a rectangle that has been dissected into nine non-overlapping squares. Given that the width and the height of the rectangle are relatively prime positive integers, find the perimeter of the rectangle.


 
                    \begin{figure}[H]
                        \centering
                        \includegraphics[width=0.5\textwidth]{Problem/AIME_2000_I_4_Problem_Figure_1.png}
                    \end{figure}
                    


 
 \textbf{Answer: }260



\textbf{Solution 1}

Call the squares' side lengths from smallest to largest $a_1,\ldots,a_9$, and let $l,w$ represent the dimensions of the rectangle.

 The picture shows that
 \begin{align*} a_1+a_2 &= a_3\\ a_1 + a_3 &= a_4\\ a_3 + a_4 &= a_5\\ a_4 + a_5 &= a_6\\ a_2 + a_3 + a_5 &= a_7\\ a_2 + a_7 &= a_8\\ a_1 + a_4 + a_6 &= a_9\\ a_6 + a_9 &= a_7 + a_8.\end{align*}


 Expressing all terms 3 to 9 in terms of $a_1$ and $a_2$ and substituting their expanded forms into the previous equation will give the expression $5a_1 = 2a_2$.

 We can guess that $a_1 = 2$. (If we started with $a_1$ odd, the resulting sides would not be integers and we would need to scale up by a factor of $2$ to make them integers; if we started with $a_1 > 2$ even, the resulting dimensions would not be relatively prime and we would need to scale down.) Then solving gives $a_9 = 36$, $a_6=25$, $a_8 = 33$, which gives us $l=61,w=69$.  These numbers are relatively prime, as desired. The perimeter is $2(61)+2(69)=\boxed{260}$.

 





\textbf{Solution 1.2 (more detail)}

We can just list the equations:
 \begin{align*} s_3 &= s_1 + s_2 \\ s_4 &= s_3 + s_1 \\ s_5 &= s_4 + s_3 \\ s_6 &= s_5 + s_4 \\ s_7 &= s_5 + s_3 + s_2 \\ s_8 &= s_7 + s_2 \\ s_9 &= s_8 + s_2 - s_1 \\ s_9 + s_8 &= s_7 + s_6 + s_5 \end{align*}
We can then write each $s_i$ in terms of $s_1$ and $s_2$ as follows
 \begin{align*} s_4 &= 2s_1 + s_2 \\ s_5 &= 3s_1 +2s_2 \\ s_6 &= 5s_1 + 3s_2 \\ s_7 &= 4s_1 + 4s_2 \\ s_8 &= 4s_1 + 5s_2 \\ s_9 &= 3s_1 + 6s_2 \\ \end{align*}
Since $s_9 + s_8 = s_7 + s_6 + s_5 \implies (3s_1 + 6s_2) + (4s_1 + 5s_2) = (4s_1 + 4s_2) + (5s_1 + 3s_2) + (3s_1 + 2s_2),$ 
 \[2s_2 = 5s_1 \implies \frac{2}{5}s_2 = s_1.\]Since the side lengths of the rectangle are relatively prime, we can see that $s_1 = 2$ and $s_2 = 5.$ Therefore, $2(2s_9 + s_6 + s_8) = 30s_1 + 40s_2 = \boxed{260}.$~peelybonehead

 





\textbf{Solution 2 Length-chasing (Angle-chasing but for side lengths)}

We set the side length of the smallest square to 1, and set the side length of square $a_4$ in the previous question to a. We do some "side length chasing" and get $4a - 4 = 2a + 5$. Solving, we get $a = 4.5$ and the side lengths are $61$ and $69$. Thus, the perimeter of the rectangle is $2(61 + 69) = \boxed{260}.$

 




	\item (\textit{AIME I 2000}) Each of two boxes contains both black and white marbles, and the total number of marbles in the two boxes is $25.$ One marble is taken out of each box randomly. The probability that both marbles are black is $\frac{27}{50},$ and the probability that both marbles are white is $\frac{m}{n},$ where $m$ and $n$ are relatively prime positive integers. What is $m + n$?


 
 \textbf{Answer: }026



\textbf{Solution 1}

If we work with the problem for a little bit, we quickly see that there is no direct combinatorics way to calculate $m/n$. The \href{https://artofproblemsolving.com/wiki/index.php?title=Principle\_of\_Inclusion-Exclusion}{Principle of Inclusion-Exclusion} still requires us to find the individual probability of each box.

 Let $a, b$ represent the number of marbles in each box, and \href{https://artofproblemsolving.com/wiki/index.php?title=Without\_loss\_of\_generality}{without loss of generality} let $a>b$. Then, $a + b = 25$, and since the $ab$ may be reduced to form $50$ on the denominator of $\frac{27}{50}$, $50|ab$. It follows that $5|a,b$, so there are 2 pairs of $a$ and $b: (20,5),(15,10)$.

 
                \begin{itemize}
                    \item \textbf{Case 1}: Then the product of the number of black marbles in each box is $54$, so the only combination that works is $18$ black in first box, and $3$ black in second. Then, $P(\text{both white}) = \frac{2}{20} \cdot \frac{2}{5} = \frac{1}{25},$ so $m + n = 26$.
                \end{itemize}


                
                \begin{itemize}
                    \item \textbf{Case 2}: The only combination that works is 9 black in both. Thus, $P(\text{both white}) = \frac{1}{10}\cdot \frac{6}{15} = \frac{1}{25}$. $m + n = 26$.
                \end{itemize}


                Thus, $m + n = \boxed{026}$.

 





\textbf{Solution 2}

Let $w_1, w_2, b_1,$ and $b_2$ represent the white and black marbles in boxes 1 and 2.

 Since there are $25$ marbles in the box:

 $w_1 + w_2 + b_1 + b_2 = 25$

 From the fact that there is a $\frac{27}{50}$ chance of drawing one black marble from each box:

 $\frac{b_1 \cdot b_2}{(b_1 + w_1)(b_2 + w_2)} = \frac{27}{50} = \frac{54}{100} = \frac{81}{150}$

 Thinking of the numerator and denominator separately, if $\frac{27}{50}$ was not a reduced fraction when calculating out the probability, then $b_1 \cdot b_2 = 27$.  Since $b_1 < 25$, this forces the variables to be $3$ and $9$ in some permutation.  Without loss of generality, let $b_1 = 3$ and $b_2 = 9$.

 The denominator becomes:$(3 + w_1)(9 + w_2) = 50$

 Since there have been $12$ black marbles used, there must be $13$ white marbles.  Substituting that in:

 $(3 + w_1)(9 + (13 - w_1)) = 50$

 $(3 + w_1)(22 - w_1) = 50$

 Since the factors of $50$ that are greater than $3$ are $5, 10, 25,$ and $50$, the quantity $3 + w_1$ must equal one of those.  However, since $w_1 < 13$, testing $2$ and $7$ for $w_1$ does not give a correct product.  Thus, $\frac{27}{50}$ must be a reduced form of the actual fraction.

 First assume that the fraction was reduced from $\frac{54}{100}$, yielding the equations $b_1\cdot b_2 = 54$ and $(b_1 + w_1)(b_2 + w_2) = 100$.Factoring $b_1 \cdot b_2 = 54$ and saying WLOG that $b_1 <  b_2 < 25$ gives $(b_1, b_2) = (3, 18)$ or $(6, 9)$.  Trying the first pair and setting the denominator equal to 100 gives:$(3 + w_1)(18 + w_2) = 100$

 
Since $w_1 + w_2 = 4$, the pairs $(w_1, w_2) = (1, 3), (2,2),$ and $(3,1)$ can be tried, since each box must contain at least one white marble.  Plugging in $w_1 = w_2 = 2$ gives the true equation $(3 + 2)(18 + 2) =100$, so the number of marbles are $(w_1, w_2, b_1, b_2) = (2, 2, 3, 18)$

 Thus, the chance of drawing 2 white marbles is $\frac{w_1 \cdot w_2 }{(w_1+ b_1)(w_2 + b_2)} = \frac{4}{100} = \frac{1}{25}$ in lowest terms, and the answer to the problem is  $1 + 25 = \boxed{026}.$

 


 For completeness, the fraction $\frac{81}{150}$ may be tested.  $150$ is the highest necessary denominator that needs to be tested, since the maximum the denominator $(w_1+ b_1)(w_2 + b_2)$ can be when the sum of all integer variables is $25$ is when the variables are $6, 6, 6,$ and $7$, in some permutation, which gives $154$.  If $b_1 \cdot b_2 = 81$, this forces $b_1 = b_2 = 9$, since all variables must be integers under $25$.  The denominator becomes $(9 + w_1)(9 + w_2) = 150$, and since there are now $25 - 18 = 7$ white marbles total, the denominator becomes $(9 + w_1)(16 - w_1) = 150$.  Testing $w_1 = 1$ gives a solution, and thus $w_2 = 6$.  The complete solution for this case is $(w_1, w_2, b_1, b_2) = (1, 6, 9, 9)$.  Although the distribution and colors of the marbles is different from the last case, the probability of drawing two white marbles is $\frac{6 \cdot 1}{ 150}$, which still simplifies to $\frac {1}{25}$.

 





\textbf{Solution 3}

We know that $\frac{27}{50} = \frac{b_1}{t_1} \cdot \frac{b_2}{t_2}$, where $b_1$ and $b_2$ are the number of black marbles in the first and the second box respectively, and $t_1$ and $t_2$ is the total number of marbles in the first and the second boxes respectively. So, $t_1 + t_2 = 25$. Then, we can realize that $\frac{27}{50} = \frac{9}{10} \cdot \frac{3}{5} = \frac{9}{10} \cdot \frac{9}{15}$, which means that having 9 black marbles out of 10 total in the first box and 9 marbles out of 15 total the second box is valid. Then there is 1 white marble in the first box and 6 in the second box. So, the probability of drawing two white marbles becomes $\frac{1}{10} \cdot \frac{6}{15} = \frac{1}{25}$. The answer is $1 + 25 = \boxed{026}$

 Note: Note that if $t_1=5, t_2=20$, and $b_1=3, b_2=18$, it also works since $\frac{b_1}{t_1} \cdot \frac{b_2}{t_2} = \frac{3}{5} \cdot \frac{18}{20} = \frac{27}{50}$, so the probability of drawing a white marble is $\frac{2}{5} \cdot \frac{2}{20} = \frac{1}{25}$. Therefore, our answer is $1+25=\boxed{026}.$

 




	\item (\textit{AIME I 2000}) For how many ordered pairs $(x,y)$ of integers is it true that $0 < x < y < 10^{6}$ and that the arithmetic mean of $x$ and $y$ is exactly $2$ more than the geometric mean of $x$ and $y$?


 
 \textbf{Answer: }997



\textbf{Solution 1}


 \begin{eqnarray*} \frac{x+y}{2} &=& \sqrt{xy} + 2\\ x+y-4 &=& 2\sqrt{xy}\\ y - 2\sqrt{xy} + x &=& 4\\ \sqrt{y} - \sqrt{x} &=& \pm 2\end{eqnarray*}

 Because $y > x$, we only consider $+2$.

 For simplicity, we can count how many valid pairs of $(\sqrt{x},\sqrt{y})$ that satisfy our equation.

 The maximum that $\sqrt{y}$ can be is $\sqrt{10^6} - 1 = 999$ because $\sqrt{y}$ must be an integer (this is because $\sqrt{y} - \sqrt{x} = 2$, an integer). Then $\sqrt{x} = 997$, and we continue this downward until $\sqrt{y} = 3$, in which case $\sqrt{x} = 1$. The number of pairs of $(\sqrt{x},\sqrt{y})$, and so $(x,y)$ is then $\boxed{997}$.

 





\textbf{Solution 2}

Let $a^2$ = $x$ and $b^2$ = $y$, where $a$ and $b$ are positive.

 Then 
 \[\frac{a^2 + b^2}{2} = \sqrt{{a^2}{b^2}} +2\]
 \[a^2 + b^2 = 2ab + 4\]
 \[(a-b)^2 = 4\]
 \[(a-b) = \pm 2\]

 This makes counting a lot easier since now we just have to find all pairs $(a,b)$ that differ by 2.

 
Because $\sqrt{10^6} = 10^3$, then we can use all positive integers less than 1000 for $a$ and $b$.

 
We know that because $x < y$, we get $a < b$.

 
We can count even and odd pairs separately to make things easier*:

 
Odd: 
 \[(1,3) , (3,5) , (5,7)  .  .  .  (997,999)\]

 
Even: 
 \[(2,4) , (4,6) , (6,8)  .  .  .  (996,998)\]

 
This makes $499$ odd pairs and $498$ even pairs, for a total of $\boxed{997}$ pairs.

 


 $*$Note: We are counting the pairs for the values of $a$ and $b$, which, when squared, translate to the pairs of $(x,y)$ we are trying to find.

 





\textbf{Solution 3}

Since the arithmetic mean is 2 more than the geometric mean, $\frac{x+y}{2} = 2 + \sqrt{xy}$. We can multiply by 2 to get $x + y = 4 + 2\sqrt{xy}$. Subtracting 4 and squaring gives 
 \[((x+y)-4)^2 = 4xy\]
 \[((x^2 + 2xy + y^2) + 16 - 2(4)(x+y)) = 4xy\]
 \[x^2 - 2xy + y^2 + 16 - 8x - 8y = 0\]

 Notice that $((x-y)-4)^2 = x^2 - 2xy + y^2 + 16 - 8x +8y$, so the problem asks for solutions of 
 \[(x-y-4)^2 = 16y\]Since the left hand side is a perfect square, and 16 is a perfect square, $y$ must also be a perfect square. Since $0 < y < (1000)^2$, $y$ must be from $1^2$ to $999^2$, giving at most 999 options for $y$.

 However if $y = 1^2$, you get $(x-5)^2 = 16$, which has solutions $x = 9$ and $x = 1$. Both of those solutions are not less than $y$, so $y$ cannot be equal to 1. If $y = 2^2 = 4$, you get $(x - 8)^2 = 64$, which has 2 solutions, $x = 16$, and $x = 0$. 16 is not less than 4, and $x$ cannot be 0, so $y$ cannot be 4. However, for all other $y$, you get exactly 1 solution for $x$, and that gives a total of $999 - 2 = \boxed{997}$ pairs.

 - asbodke

 


 





\textbf{Solution 4 (Similar to Solution 3)}

Rearranging our conditions to 

 
 \[x^2-2xy+y^2+16-8x-8y=0 \implies\]
 \[(y-x)^2=8(x+y-2).\]

 Thus, $4|y-x.$

 Now, let $y = 4k+x.$ Plugging this back into our expression, we get

 
 \[(k-1)^2=x.\]

 There, a unique value of $x, y$ is formed for every value of $k$. However, we must have 

 
 \[y<10^6 \implies (k+1)^2< 10^6-1\]

 and 

 
 \[x=(k-1)^2+1>0.\]

 Therefore, there are only $\boxed{997}$ pairs of $(x,y).$

 Solution by Williamgolly

 





\textbf{Solution 5}

First we see that our condition is $\frac{x+y}{2} = 2 + \sqrt{xy}$. Then we can see that $x+y = 4 + 2\sqrt{xy}$. From trying a simple example to figure out conditions for $x,y$, we want to find $x-y$ so we can isolate for $x$. From doing the example we can note that we can square both sides and subtract $4xy$: $(x-y)^2 = 16 + 16\sqrt{xy} \implies x-y = -2( \sqrt{1+\sqrt{xy}})$ (note it is negative because $y > x$. Clearly the square root must be an integer, so now let $\sqrt{xy} = a^2-1$. Thus $x-y = -2a$. Thus $x = 2 + \sqrt{xy} - a = 2 + a^2 - 1 -2a$. We can then find $y$, and use the quadratic formula on $x,y$ to ensure they are $>0$ and $<10^6$ respectively. Thus we get that $y$ can go up to 999 and $x$ can go down to $3$, leaving $997$ possibilities for $x,y$.

 




	\item (\textit{AIME I 2000}) Suppose that $x,$ $y,$ and $z$ are three positive numbers that satisfy the equations $xyz = 1,$ $x + \frac {1}{z} = 5,$ and $y + \frac {1}{x} = 29.$ Then $z + \frac {1}{y} = \frac {m}{n},$ where $m$ and $n$ are relatively prime positive integers. Find $m + n$.


 
 \textbf{Answer: }005



\textbf{Solution 1}

We can rewrite $xyz=1$ as $\frac{1}{z}=xy$.

 Substituting into one of the given equations, we have 
 \[x+xy=5\]
 \[x(1+y)=5\]
 \[\frac{1}{x}=\frac{1+y}{5}.\]

 We can substitute back into $y+\frac{1}{x}=29$ to obtain
 \[y+\frac{1+y}{5}=29\]
 \[5y+1+y=145\]
 \[y=24.\]

 We can then substitute once again to get
 \[x=\frac15\]
 \[z=\frac{5}{24}.\]Thus, $z+\frac1y=\frac{5}{24}+\frac{1}{24}=\frac{1}{4}$, so $m+n=\boxed{005}$.

 





\textbf{Solution 2}

Let $r = \frac{m}{n} = z + \frac {1}{y}$.

 
 \begin{align*} (5)(29)(r)&=\left(x + \frac {1}{z}\right)\left(y + \frac {1}{x}\right)\left(z + \frac {1}{y}\right)\\ &=xyz + \frac{xy}{y} + \frac{xz}{x} + \frac{yz}{z} + \frac{x}{xy} + \frac{y}{yz} + \frac{z}{xz} + \frac{1}{xyz}\\ &=1 + x + z + y + \frac{1}{y} + \frac{1}{z} + \frac{1}{x} + \frac{1}{1}\\ &=2 + \left(x + \frac {1}{z}\right) + \left(y + \frac {1}{x}\right) + \left(z + \frac {1}{y}\right)\\ &=2 + 5 + 29 + r\\ &=36 + r \end{align*}


 Thus $145r = 36+r \Rightarrow 144r = 36 \Rightarrow r = \frac{36}{144} = \frac{1}{4}$. So $m + n = 1 + 4 = \boxed{5}$.

 





\textbf{Solution 3}

Since $x+\frac{1}{z}=5, 1=z(5-x)=xyz$, so $5-x=xy$. Also, $y=29-\frac{1}{x}$ by the second equation. Substitution gives $x=\frac{1}{5}$, $y=24$, and $z=\frac{5}{24}$, so the answer is $4+1$ which is equal to $5$.

 





\textbf{Solution 4}

(Hybrid between 1/2)

 Because $xyz = 1, \hspace{0.15cm} \frac{1}{x} = yz, \hspace{0.15cm} \frac{1}{y} = xz,$ and $\hspace{0.05cm}\frac{1}{z} = xy$. Substituting and factoring, we get $x(y+1) = 5$, $\hspace{0.15cm}y(z+1) = 29$, and $\hspace{0.05cm}z(x+1) = k$. Multiplying them all together, we get, $xyz(x+1)(y+1)(z+1) = 145k$, but $xyz$ is $1$, and by the Identity property of multiplication, we can take it out. So, in the end, we get $(x+1)(y+1)(z+1) = 145k$. And, we can expand this to get $xyz+xy+yz+xz+x+y+z+1 = 145k$, and if we make a substitution for $xyz$, and rearrange the terms, we get $xy+yz+xz+x+y+z = 145k-2$ This will be important.

 
Now, lets add the 3 equations $x(y+1) = 5, \hspace{0.15cm}y(z+1) = 29$, and $\hspace{0.05cm}z(x+1) = k$. We use the expand the Left hand sides, then, we add the equations to get $xy+yz+xz+x+y+z = k+34$ Notice that the LHS of this equation matches the LHS equation that I said was important. So, the RHS of both equations are equal, and thus $145k-2 = k+34$ We move all constant terms to the right, and all linear terms to the left, to get $144k = 36$, so $k = \frac{1}{4}$ which gives an answer of $1+4 = \boxed{005}$

 -AlexLikeMath

 





\textbf{Solution 5}

Get rid of the denominators in the second and third equations to get $xz-5z=-1$ and $xy-29x=-1$. Then, since $xyz=1$, we have $\frac 1y-5z=-1$ and $\frac 1z-29x=-1$. Then, since we know that $\frac 1z+x=5$, we can subtract these two equations to get that $30x=6\implies x=5$. The result follows that $z=\frac 5{24}$ and $y=24$, so $z+\frac 1y=\frac 1{24}+\frac 5{24}=\frac 14$, and the requested answer is $1+4=\boxed{005}.$

 





\textbf{Solution 6}

Rewrite the equations in terms of x. 

 $x+\frac{1}{z}=5$ becomes $z=\frac{1}{x+5}$. 

 $y+\frac{1}{x}=29$ becomes $y=29-\frac{1}{x}$

 Now express $xyz=1$ in terms of x. 

 $\frac{1}{5-x}\cdot(29-\frac{1}{x})\cdot x=1$.

 This evaluates to $29x-1=5-x$, giving us $x=\frac{1}{5}$. We can now plug x into the other equations to get $y=24$ and $z=\frac{5}{24}$. 

 Therefore, $z+\frac{1}{y}=\frac{5}{24}+\frac{1}{24}=\frac{6}{24}=\frac{1}{4}$. 

 $1+4=\boxed{005}$, and we are done.~MC413551

 





\textbf{Solution 7 (note: identical to Solution 2)}

First, we have been given the value of $xyz$, so we should probably figure out a way to use that. Before that, we replace $\dfrac{m}{n}$, because why deal with $2$ variables when you don't have to. We multiply these $3$ equations to get: 
 \[(x+\dfrac{1}{z})(y+\dfrac{1}{x})(z+\dfrac{1}{y}) \Longrightarrow xyz+x+y+z+\dfrac{1}{x}+\dfrac{1}{y}+\dfrac{1}{z} = 145a\]. Plugging the value in for $xyz$, we get: 
 \[2+x+y+z+\dfrac{1}{x}+\dfrac{1}{y}+\dfrac{1}{z}\]. Now, we need a way to get these other random terms out of there. We add the three equations to get 
 \[x+y+z+\dfrac{1}{x}+\dfrac{1}{y}+\dfrac{1}{z} = 34+a\]. Plugging this back into the equation we found we get: 
 \[36+a = 145a \Longrightarrow a = \dfrac{1}{4} \Longrightarrow \boxed{005}\]

 -jb2015007

 





\textbf{Solution 8}

Let $x=\frac{a}{b},y=\frac{b}{c},$ and $z=\frac{c}{a}.$ From $x+\frac{1}{z}=5,$ we get $\frac{a}{b}+\frac{a}{c}=5\Rightarrow ac+ab=5bc.$ Similarly, from $y+\frac{1}{x}=29,ab+bc=29ac.$ We wish to find $\frac{c}{a}+\frac{c}{b}=\frac{ac+bc}{ab}.$ Set $ab=1.$ Then $ac+1=5bc$ and $bc+1=29ac.$ Solving these equations, we get $bc=\frac{5}{24}$ and $ac=\frac{1}{24}$ and thus we get $\frac{5}{24}+\frac{1}{24}=\frac{1}{4}\Rightarrow m=1,n=4$ and thus our answer is $\boxed{005}.$~rice\_farmer

 




	\item (\textit{AIME I 2000}) A container in the shape of a right circular cone is 12 inches tall and its base has a 5-inch radius. The liquid that is sealed inside is 9 inches deep when the cone is held with its point down and its base horizontal. When the liquid is held with its point up and its base horizontal, the height of the liquid is $m - n\sqrt [3]{p},$ where $m,$ $n,$ and $p$ are positive integers and $p$ is not divisible by the cube of any prime number. Find $m + n + p$.


 
 \textbf{Answer: }052



\textbf{Solution}


                    \begin{figure}[H]
                        \centering
                        \includegraphics[width=0.5\textwidth]{Solution/AIME_2000_I_8_Solution_Figure_1.png}
                    \end{figure}
                    

 





\textbf{Solution 1}

The scale factor is uniform in all dimensions, so the volume of the liquid is $\left(\frac{3}{4}\right)^{3}$ of the container. The remaining section of the volume is $\frac{1-\left(\frac{3}{4}\right)^{3}}{1}$ of the volume, and therefore $\frac{\left(1-\left(\frac{3}{4}\right)^{3}\right)^{1/3}}{1}$ of the height when the vertex is at the top.

 So, the liquid occupies $\frac{1-\left(1-\left(\frac{3}{4}\right)^{3}\right)^{1/3}}{1}$ of the height, or $12-12\left(1-\left(\frac{3}{4}\right)^{3}\right)^{1/3}=12-3\left(37^{1/3}\right)$. Thus $m+n+p=\boxed{052}$.

 





\textbf{Solution 2}

(Computational) The volume of a cone can be found by $V = \frac{\pi}{3}r^2h$. In the second container, if we let $h',r'$ represent the height, radius (respectively) of the air (so $12 -h'$ is the height of the liquid), then the volume of the liquid can be found by $\frac{\pi}{3}r^2h - \frac{\pi}{3}(r')^2h'$. 

 By \href{https://artofproblemsolving.com/wiki/index.php?title=Similar\_triangle&action=edit&redlink=1}{similar triangles}, we find that the dimensions of the liquid in the first cone to the entire cone is $\frac{3}{4}$, and that $r' = \frac{rh'}{h}$; equating,

 
 \begin{align*}\frac{\pi}{3}\left(\frac{3}{4}r\right)^2 \left(\frac{3}{4}h\right) &= \frac{\pi}{3}\left(r^2h - \left(\frac{rh'}{h}\right)^2h'\right)\\ \frac{37}{64}r^2h &= \frac{r^2}{h^2}(h')^3 \\ h' &= \sqrt[3]{\frac{37}{64} \cdot 12^3} = 3\sqrt[3]{37}\end{align*}


 Thus the answer is $12 - h' = 12-3\sqrt[3]{37}$, and $m+n+p=\boxed{052}$.

 





\textbf{Solution 3}

From the formula $V=\frac{\pi r^2h}{3}$, we can find that the volume of the container is $100\pi$. The cone formed by the liquid is similar to the original, but scaled down by $\frac{3}{4}$ in all directions, so its volume is $100\pi*\frac{27}{64}=\frac{675\pi}{16}$. The volume of the air in the container is the volume of the container minus the volume of the liquid, which is $\frac{925\pi}{16}$, which is $\frac{37}{64}$ of the volume of the container. When the point faces upwards, the air forms a cone at the top of the container. This cone must have $\sqrt[3]{\frac{37}{64}}=\frac{\sqrt[3]{37}}{4}$ of the height of the container. This means that the height of the liquid is $12\left(1-\frac{\sqrt[3]{37}}{4}\right)=12-3\sqrt[3]{37}$ inches, so our answer is $\boxed{052}$. Solution by Zeroman

 





\textbf{Solution 4}

We find that the volume of the cone is $100\pi$. 

 The volume of the cone with height 9 is $\frac{675}{16}\pi$. 

 The difference between the two volumes is $\frac{925}{16}\pi$. Note that this is the volume of the cone essentially 'on top of' the frustum described in the problem when the liquid is held with the base horizontal. 

 We can express the volume as $x^2\pi\cdot\frac{12}{5}x\cdot\frac{1}{3}$, where x is the radius of the cone. 

 Solving this equation, we get $x=\frac{5\sqrt[3]{37}}{4}$. The height of this cone is $\frac{12}{5}$ of the radius. 

 Then we subtract the value from the height, 12, to get our answer: $12-3\sqrt[3]{37}$. 

 Therefore, our answer is $12+3+37=\boxed{52}$. ~MC413551

 




	\item (\textit{AIME I 2000}) The system of equations

 \begin{eqnarray*}\log_{10}(2000xy) - (\log_{10}x)(\log_{10}y) & = & 4 \\ \log_{10}(2yz) - (\log_{10}y)(\log_{10}z) & = & 1 \\ \log_{10}(zx) - (\log_{10}z)(\log_{10}x) & = & 0 \\ \end{eqnarray*}


 has two solutions $(x_{1},y_{1},z_{1})$ and $(x_{2},y_{2},z_{2})$. Find $y_{1} + y_{2}$.


 
 \textbf{Answer: }025



\textbf{Solution}

Since $\log ab = \log a + \log b$, we can reduce the equations to a more recognizable form:

 
 \begin{eqnarray*} -\log x \log y + \log x + \log y - 1 &=& 3 - \log 2000\\ -\log y \log z + \log y + \log z - 1 &=& - \log 2\\ -\log x \log z + \log x + \log z - 1 &=& -1\\ \end{eqnarray*}

 Let $a,b,c$ be $\log x, \log y, \log z$ respectively. Using \href{https://artofproblemsolving.com/wiki/index.php?title=Simon\%27s\_Favorite\_Factoring\_Trick}{SFFT}, the above equations become (*)

 
 \begin{eqnarray*}(a - 1)(b - 1) &=& \log 2 \\ (b-1)(c-1) &=& \log 2 \\ (a-1)(c-1) &=& 1  \end{eqnarray*}

 Small note from different author: $-(3 - \log 2000) = \log 2000 - 3 = \log 2000 - \log 1000 = \log 2.$

 From here, multiplying the three equations gives 

 
 \begin{eqnarray*}(a-1)^2(b-1)^2(c-1)^2 &=& (\log 2)^2\\ (a-1)(b-1)(c-1) &=& \pm\log 2\end{eqnarray*}

 Dividing the third equation of (*) from this equation, $b-1 = \log y - 1 = \pm\log 2 \Longrightarrow \log y = \pm \log 2 + 1$. (Note from different author if you are confused on this step: if $\pm$ is positive then $\log y = \log 2 + 1 = \log 2 + \log 10 = \log 20,$ so $y=20.$ if $\pm$ is negative then $\log y = 1 - \log 2 = \log 10 - \log 2 = \log 5,$ so $y=5.$) This gives $y_1 = 20, y_2 = 5$, and the answer is $y_1 + y_2 = \boxed{025}$.

 





\textbf{Solution 2}

Subtracting the second equation from the first equation yields
 \begin{align*} \log 2000xy-\log 2yz-((\log x)(\log y)-(\log y)(\log z)) &= 3 \\ \log\frac{2000xy}{2yz}-\log y(\log x-\log z) &= 3 \\ \log1000+\log\frac{x}{z}-\log y(\log\frac{x}{z}) &= 3 \\ 3+\log\frac{x}{z}-\log y(\log\frac{x}{z}) &= 3 \\ \log\frac{x}{z}(1-\log y) &= 0 \\ \end{align*}
If $1-\log y=0$ then $y=10$. Substituting into the first equation yields $\log20000=4$ which is not possible.

 If $\log\frac{x}{z}=0$ then $\frac{x}{z}=1\Longrightarrow x=z$. Substituting into the third equation gets
 \begin{align*} \log x^2-(\log x)(\log x) &= 0 \\ \log x^2-\log x^x &= 0 \\ \log x^{2-x} &= 0 \\ x^{2-x} &= 1 \\ \end{align*}
Thus either $x=1$ or $2-x=0\Longrightarrow x=2$. (Note that here $x\neq-1$ since logarithm isn't defined for negative number.)

 Substituting $x=1$ and $x=2$ into the first equation will obtain $y=5$ and $y=20$, respectively. Thus $y_1+y_2=\boxed{025}$.

 





\textbf{Solution 3}

Let $a = \log x$, $b = \log y$ and $c = \log z$. Then the given equations become: 

 
 \begin{align*} \log 2 + a + b - ab = 1 \\ \log 2 + b + c - bc = 1 \\ a+c = ac \\ \end{align*}


 Equating the first and second equations, solving, and factoring, we get $a(1-b) = c(1-b) \implies{a = c}$. Plugging this result into the third equation, we get $c = 0$ or $2$. Substituting each of these values of $c$ into the second equation, we get $b = 1 - \log 2$ and $b = 1 + \log 2$. Substituting backwards from our original substitution, we get $y = 5$ and $y = 20$, respectively, so our answer is $\boxed{025}$.

 





\textbf{Solution 4}

All logs are base 10 by convention. Rearrange the given statements:

 
 \begin{align*} \log 2000 + \log x + \log y - \log x \log y = 4, \quad \text{which becomes} \quad \log x + \log y = \log x \log y + \log 5. \\ \log 2 + \log y + \log z - \log y \log z = 1, \quad \text{which becomes} \quad \log y + \log z = \log y \log z + \log 5. \\ \log z + \log x - \log z \log x = 0, \quad \text{which becomes} \quad \log x + \log z = \log z \log x. \\ \end{align*}


 Subtract the first two equations to obtain $(\log x - \log z) = \log y (\log x - \log z).$

 This must mean \(\log x = \log z\) because otherwise \(\log y = 1\) turns the first equation into \(\log y = \log 5\) which is self-contradictory.

 With \(\log x = \log z\) we know each value satisfies \(2a = a^2\), so they are both \(0\) or both \(2\). Finally we arrive at our two solutions, where 0 gives us \(0 + \log y = 0 + \log 5\), and \(y = 5\), and 2 gives us \(2 + \log y = 2 \log y + \log 5\), and \(\log y = 2 - \log 5 = \log 20\), so \(y = 20\). Similar to above we arrive at $\boxed{025}$.

 




	\item (\textit{AIME I 2000}) A sequence of numbers $x_{1},x_{2},x_{3},\ldots,x_{100}$ has the property that, for every integer $k$ between $1$ and $100,$ inclusive, the number $x_{k}$ is $k$ less than the sum of the other $99$ numbers. Given that $x_{50} = \frac{m}{n}$, where $m$ and $n$ are relatively prime positive integers, find $m + n$.


 
 \textbf{Answer: }173



\textbf{Solution}

Let the sum of all of the terms in the sequence be $\mathbb{S}$. Then for each integer $k$, $x_k = \mathbb{S}-x_k-k \Longrightarrow \mathbb{S} - 2x_k = k$. Summing this up for all $k$ from $1, 2, \ldots, 100$,

 
 \begin{align*}100\mathbb{S}-2(x_1 + x_2 + \cdots + x_{100}) &= 1 + 2 + \cdots + 100\\ 100\mathbb{S} - 2\mathbb{S} &= \frac{100 \cdot 101}{2} = 5050\\ \mathbb{S}&=\frac{2525}{49}\end{align*}


 Now, substituting for $x_{50}$, we get $2x_{50}=\frac{2525}{49}-50=\frac{75}{49} \Longrightarrow x_{50}=\frac{75}{98}$, and the answer is $75+98=\boxed{173}$.

 





\textbf{Solution 2}

Consider $x_k$ and $x_{k+1}$. Let $S$ be the sum of the rest 98 terms. Then $x_k+k=S+x_{k+1}$ and $x_{k+1}+(k+1)=S+x_k.$ Eliminating $S$ we have $x_{k+1}-x_k=-\dfrac{1}{2}.$ So the sequence is arithmetic with common difference $-\dfrac{1}{2}.$ 

 In terms of $x_{50},$ the sequence is $x_{50}+\dfrac{49}{2}, x_{50}+\dfrac{48}{2},\cdots,x_{50}+\dfrac{1}{2}, x_{50}, x_{50}-\dfrac{1}{2}, \cdots, x_{50}-\dfrac{49}{2}, x_{50}-\dfrac{50}{2}.$ Therefore, $x_{50}+50=99x_{50}-\dfrac{50}{2}$.

 Solving, we get $x_{50}=\dfrac{75}{98}.$ The answer is $75+98=\boxed{173}.$

 - JZ

 - edited by erinb28lms

 





\textbf{Solution 3 (Sum of equations)}

Like Solution 1, let the sum of all of the terms in this sequence be $\mathbb{S}$. By definition:

 
 \[x_1 + 1 = x_2+x_3+x_4+...+x_{100}\]
 \[x_2 + 2 = x_1+x_3+x_4+...+x_{100}\]
 \[x_3 + 3 = x_1+x_2+x_4+...+x_{100}\]
 \[...\]
 \[x_{99} + 99 = x_1+x_2+x_3+...+x_{98}+x_{100}\]
 \[x_{100} + 100 = x_1+x_2+x_3+...+x_{98}+x_{99}\].

 Adding up all of these equations yields:

 
 \[\mathbb{S} + TR(100) = 99\mathbb{S}\]

 Here $TR(100)$ represents the $100$th triangular number, which is $5050$. Solving for $\mathbb{S}$ yields:

 
 \[\mathbb{S} = \frac{2525}{49}\].

 $\mathbb{S}$ can also be written as $x_{50} + (x_{50} + 50)$. Solving for $x_{50}$,

 
 \[2x_{50} = \frac{2525-2450}{49}\]
 \[2x_{50} = \frac{75}{49}\]
 \[x_{50} = \frac{75}{98}\]

 The requested sum is therefore $75+98 = \boxed{173}$.

 





\textbf{Solution 4 (Clever Substitution)}

Due to the problem's unique setup, we can actually represent every $x_k$ in terms of its sum, call it $S$. For any $x_k$,
 \[x_k = \frac{S}{2} - \frac{k}{2}\]We can show that this works because the sum of the other $99$ terms is now equal to $\frac{S}{2} + \frac{k}{2}$, and the difference is indeed $k$. The motivation to find that we need an $\frac{S}{2}$ term lies in that the difference between $x_k$ and the $99$ other terms is an integer, and not an expression in terms of $S$, and allowing $x_k$ to be in terms of $\frac{S}{2}$ is the only way to achieve this.

 Now, finding the numerical value of $S$ solves the problem. Summing all terms from $x_1$ to $x_{100}$, we find
 \[S = \sum_{k=1}^{100} x_k = 100\frac{S}{2} - \sum_{k=1}^{100} \frac{k}{2} = 50S - \frac{100\cdot\frac{1 + 100}{2}}{2} = 50S - 2525\]

 We obtain $S = \frac{2525}{49}$, and we now only need to find $x_{50}$, which is represented as $\frac{S}{2} - \frac{50}{2}$. Substituting, 
 \[x_{50} = \frac{2525}{49\cdot2} - 25 = \frac{25\cdot101 - 25\cdot98}{98} = \frac{75}{98}\]Our desired answer is thus $75 + 98 = \boxed{173}$

 




	\item (\textit{AIME I 2000}) Let $S$ be the sum of all numbers of the form $\frac{a}{b}$, where $a$ and $b$ are relatively prime positive divisors of $1000.$ What is the greatest integer that does not exceed $\frac{S}{10}$?


 
 \textbf{Answer: }248



\textbf{Solution 1}

Since all divisors of $1000 = 2^35^3$ can be written in the form of $2^{m}5^{n}$, it follows that $\frac{a}{b}$ can also be expressed in the form of $2^{x}5^{y}$, where $-3 \le x,y \le 3$. Thus every number in the form of $a/b$ will be expressed one time in the product

 
 \[(2^{-3} + 2^{-2} + 2^{-1} + 2^{0} + 2^{1} + 2^2 + 2^3)(5^{-3} + 5^{-2} +5^{-1} + 5^{0} + 5^{1} + 5^2 + 5^3)\]

 Using the formula for a \href{https://artofproblemsolving.com/wiki/index.php?title=Geometric\_series}{geometric series}, this reduces to $S = \frac{2^{-3}(2^7 - 1)}{2-1} \cdot \frac{5^{-3}(5^{7} - 1)}{5-1} = \frac{127 \cdot 78124}{4000} = 2480 + \frac{437}{1000}$, and $\left\lfloor \frac{S}{10} \right\rfloor = \boxed{248}$.

 





\textbf{Solution 2}

Essentially, the problem asks us to compute 
 \[\sum_{a=-3}^3 \sum_{b=-3}^3 \frac{2^a}{5^b}\] which is pretty easy: 
 \[\sum_{a=-3}^3 \sum_{b=-3}^3 \frac{2^a}{5^b} = \sum_{a=-3}^3 2^a \sum_{b=-3}^3 \frac{1}{5^b} = \sum_{a=-3}^3 2^a 5^{3}\bigg( \frac{1-5^{-7}}{1-\frac{1}{5}}    \bigg) = 5^{3}\bigg( \frac{1-5^{-7}}{1-\frac{1}{5}} \bigg) \sum_{a=-3}^3 2^a = 5^{3}\bigg( \frac{1-5^{-7}}{1-\frac{1}{5}} \bigg)2^{-3} \bigg( \frac{1-2^7}{1-2}  \bigg) = 2480 + \frac{437}{1000}\] so our answer is $\left\lfloor \frac{2480 + \frac{437}{1000}}{10} \right\rfloor = \boxed{248}$.

 


 





\textbf{Solution 3}

The sum is equivalent to $\sum_{i | 10^6}^{} \frac{i}{1000}$Therefore, it's the sum of the factors of $10^6$ divided by $1000$. The sum is $\frac{127 \times 19531}{1000}$ by the sum of factors formula. The answer is therefore $\boxed{248}$ after some computation.- whatRthose

 





\textbf{Solution 4 (head on)}

We can organize the fractions and reduce them in quantities to reach our answer. First, separate the fractions with coprime parts into those that are combinations of powers of 2 and 5, and those that are a combination of a 1 and another divisor. 

 To begin with the first list, list powers of 2 and 5 from 0 to 3. In this specific case I find it easier to augment every denominator to 1000 and then divide by 1000. To do this, write the corresponding divisor under each power. e.g. 2 - 500, 4 - 250, 5 - 200, etc. Call this the "partner" of any divisor.

 Now we now the amount to multiply the numerator if given number is in the denominator. Now we simply combine and reduce these groups. If the powers of 2 are on the denominator, then every power of five will be multiplied by the partner of the power of 2. Essentially, all we have to do is a large scale distributive property application. There is nothing complicated about this except to be careful. 

 Add all powers of 2: 15

 Add their partners: 1875

 Add all powers of 5: 156

 Add their partners: 1248

 
Then, follow this formula: (sum of powers * sum opposite power's partners)+(sum of powers * sum opposite power's partners)

 Or, $156*1875+15*1248$ $=311220$

 Now, divide by 1000 to compensate for the denominator. $311.22$

 
Finally, we have to calculate the other list of fractions with 1 and another divisor. e.g. 1 - 250, 1 - 20 etc. (these all count)This time we need to list all divisors of 1000, including 1. Remove all powers of 2 or 5, because we already included those in the other list. Now, notice there are two cases. 1: 1 is in the denominator, making the fraction an integer. 2: 1 is in the numerator.

 Adding all the integers in the first case gives us 2169. The second case can actually be discarded, but still can be found. Realize that if we include the powers of 2 and 5, then the second case is (sum of all divisors)/1000. Remove all partners of powers of 2 and 5, and we'll get exactly 217/1000, or 0.22.

 
Finally, add all the numbers together: $311.22 +  2169 + 0.22 = 2480.44$

 And divide by 10: $248.044$

 After an odyssey of bashing: $\boxed{248}$

 
-jackshi2006

 





\textbf{Solution 5 (less brain)}

Since $1$ is relatively prime to everything, including itself, we start with $b=1$. The numerators will be all factors of 1000, and the sum of the factors of 1000 is $(1+2+4+8)(1+5+25+125) = 15*156=2340$.

 
If the denominator is in the form $2^k$, then the numerator must be in the form $5^j$. The sum of the possible numerators is $1+5+25+125 = 156$, so the sum of all such fractions with denominator $2^k$ is $\frac{156}{2}+\frac{156}{4}+\frac{156}{8} = 136.5$

 
If the denominator is in the form $5^k$, then the numerator must be in the form $2^j$. The sum of the possible numerators is $1+2+4+8 =15$, so the sum of all such fractions with denominator $5^k$ is $\frac{15}{5}+\frac{15}{25}+\frac{15}{125}$, which is around $3.6$.

 
If the denominator is in the form $2^k5^j$, where $k\ge 1$ and $j\ge 1$, the numerator must be $1$. We can get the sum of all such fractions from $(\frac{1}{2}+\frac{1}{4}+\frac{1}{8})(\frac{1}{5}+\frac{1}{25}+\frac{1}{125}) = \frac{217}{1000}$.

 
Adding, we get $2480$, with a bit extra, so our answer is $\boxed{248}$.

 




	\item (\textit{AIME I 2000}) Given a function $f$ for which

 \[f(x) = f(398 - x) = f(2158 - x) = f(3214 - x)\]
holds for all real $x,$ what is the largest number of different values that can appear in the list $f(0),f(1),f(2),\ldots,f(999)$?


 
 \textbf{Answer: }177



\textbf{Solution}


 \begin{align*}f(2158 - x) = f(x) &= f(3214 - (2158 - x)) &= f(1056 + x)\\ f(398 - x) = f(x) &= f(2158 - (398 - x)) &= f(1760 + x)\end{align*}


 Since $\mathrm{gcd}(1056, 1760) = 352$ we can conclude that (by the \href{https://artofproblemsolving.com/wiki/index.php?title=Euclidean\_algorithm}{Euclidean algorithm})

 
 \[f(x) = f(352 + x)\]

 So we need only to consider one period $f(0), f(1), ... f(351)$, which can have at most $352$ distinct values which determine the value of $f(x)$ at all other integers.  

 But we also know that $f(x) = f(46 - x) = f(398 - x)$, so the values $x = 24, 25, ... 46$ and $x = 200, 201, ... 351$ are repeated.  This gives a total of

 
 \[352 - (46 - 24 + 1) - (351 - 200 + 1) = \boxed{ 177 }\]

 distinct values.

 To show that it is possible to have $f(23), f(24), \ldots, f(199)$ distinct, we try to find a function which fulfills the given conditions. A bit of trial and error would lead to the \href{https://artofproblemsolving.com/wiki/index.php?title=Cosine}{cosine} function: $f(x) = \cos \left(\frac{360}{352}(x-23)\right)$ (in degrees).

 





\textbf{Solution 2}

This gives the intuition the first solution uses to solve the problem.

 One can imagine that there must be multiple lines of symmetry for the function $f(x)$, as if a function can be expressed with $f(x)=f(a-x)$ it must be symmetric against line $x=\frac{a}{2}$. Try this yourself by graphing a polynomial $f(x)$, then graphing $f(n-x)$. If $f(x)=f(n-x)$, their point of intersection at $x=\frac{n}{2}$ must contain a line of symmetry.

 For this particular function $f(x)$, it has $3$ lines of symmetry already given: $x=199$, $x=1079$, and $x=1607$. Now imagine these lines of symmetry folding over each other to form new lines of symmetry, because $x=199$ should have another corresponding line of symmetry when being reflected over the lines $x=1079$ and $x=1607$ due to it being, well, symmetric. Doing so, we find that there will be two new lines of symmetry generated by folding $x=199$ over the other 2 lines of symmetry, namely, $x=1079+(1079-199)=1959$ and $x=1607+(1607-199)=3015$. Now, if we fold $x=1079$ over the other 2 lines of symmetry, we will find another 2 lines of symmetry. We can even fold lines of symmetry not originally given by the problem, such as $x=3015$ over $x=1959$ to form even more lines of symmetry. This will result in infinite lines of symmetry, which tells us that this function is periodic.

 As we have just proven that $f(x)$ has a finite amount of values, now, we need to find what the half period it is, and not a whole period because in periodic functions such as $\cos(x)$, the $y$ values will repeat every half period and not a full period. To do this, we go back to the 3 lines of symmetry given by the problem, and we fold $x=1607$ over $x=1079$ to obtain $x=1079-(1607-1079)=551$. We now fold $x=551$ over $x=199$ to find $x=199-(551-199)=-153$. Now we fold $x=-153$ over $x=551$ to get $x=551+(551-(-153))=1255$. Now we fold $x=1255$ over $x=1079$ to find $x=903$, then we fold $x=1079$ over $x=903$, then fold $x=903$ over that line again, until suddenly, we find that after doing a few folds, there will be a line that coincides with $x=199$. Now, if you visualize a graph of all the lines of symmetry, you will find that from $x=199$ to $x=1079$, there will lines of symmetry evenly spaced by $176$ units each. If fold $x=1255$ and $x=1079$ in the other direction, you will find that it will coincide with all other named lines of symmetry and extend to infinity, with one line of symmetry every $176$ units. This tells us that the half period of $f(x)$ is $176$ units, which means that every $176$ units, the y-values will repeat. Now, all we have to do is put in the answer $\fbox{177}$.

 


 Why $177$, you ask, and not $176$? 

 Imagine a square composed of unit squares. You can find the amount of unit squares on the perimeter of the large square by taking the side length of the larger square, subtracting one, then multiplying by 4. This is not the same as finding the perimeter because if you simply found the perimeter, you would have overcounted the 4 squares on the corners. Using the same intuition, in our periodic function $f(x)$, we see that every $176$ units it repeats. Every $176$ units, it will have counted not the total amount of values, but the total amount of values minus one. That's why we need to add 1 to our answer.

 


 -WhatdoHumanitariansEat

 




	\item (\textit{AIME I 2000}) In the middle of a vast prairie, a firetruck is stationed at the intersection of two perpendicular straight highways. The truck travels at $50$ miles per hour along the highways and at $14$ miles per hour across the prairie. Consider the set of points that can be reached by the firetruck within six minutes. The area of this region is $\frac{m}{n}$ square miles, where $m$ and $n$ are relatively prime positive integers. Find $m + n$.


 
 \textbf{Answer: }731



\textbf{Solution}

Let the intersection of the highways be at the origin $O$, and let the highways be the x and y axes. We consider the case where the truck moves in the positive x direction. 

 After going $x$ miles, $t=\frac{d}{r}=\frac{x}{50}$ hours has passed. If the truck leaves the highway it can travel for at most $t=\frac{1}{10}-\frac{x}{50}$ hours, or $d=rt=14t=1.4-\frac{7x}{25}$ miles. It can end up anywhere off the highway in a circle with this radius centered at $(x,0)$. All these circles are \href{https://artofproblemsolving.com/wiki/index.php?title=Homothety}{homothetic} with respect to a center at $(5,0)$. 

 
                    \begin{figure}[H]
                        \centering
                        \includegraphics[width=0.5\textwidth]{Solution/AIME_2000_I_13_Solution_Figure_1.png}
                    \end{figure}
                         
                    \begin{figure}[H]
                        \centering
                        \includegraphics[width=0.5\textwidth]{Solution/AIME_2000_I_13_Solution_Figure_2.png}
                    \end{figure}
                    Now consider the circle at $(0,0)$. Draw a line tangent to it at $A$ and passing through $B (5,0)$. By the Pythagorean Theorem $AB^2+AO^2=OB^2 \Longrightarrow AB=\sqrt{OB^2-AO^2}=\sqrt{5^2-1.4^2}=\frac{24}{5}$. Then $\tan(\angle ABO)=\frac{OA}{AB}=\frac{7}{24}$, so the \href{https://artofproblemsolving.com/wiki/index.php?title=Slope}{slope} of line $AB$ is $\frac{-7}{24}$. Since it passes through $(5,0)$ its equation is $y=\frac{-7}{24}(x-5)$. 

 This line and the x and y axis bound the region the truck can go if it moves in the positive x direction. Similarly, the line $y=5-\frac{24}{7}x$ bounds the region the truck can go if it moves in positive y direction. The intersection of these two lines is $\left(\frac{35}{31},\frac{35}{31}\right)$. 

 
First method: The bounded region in Quadrant I is made up of two triangles each identical to the triangle with vertices $A$, $B$, and the intersection point $\left(\frac{35}{31},\frac{35}{31}\right)$. Thus, one height of this triangle is $\frac{35}{31}$ (y-coordinate of intersection point) and the corresponding base is $5$ (the entire x-length of the triangle). Then the total area in Quadrant I is $2(\frac{1}{2}\cdot\frac{35}{31}\cdot5)$. By symmetry, the regions in the other quadrants are the same, so the area of the whole region is $20x=\frac{700}{31}$ so the answer is $700+31=\boxed{731}$.

 
Second method: The bounded region in Quadrant I is made up of a square and two triangles. $A=x^2+x(5-x)=5x$. By symmetry, the regions in the other quadrants are the same, so the area of the whole region is $20x=\frac{700}{31}$ so the answer is $700+31=\boxed{731}$.

 \textbf{Intuition} Here's a bit of intuition behind the solution. If I wanted to travel to some point on the prairie, normally I would go the straight line distance. Right? However, here I would choose to travel on the highway the greatest distance, then drive the rest on the prairie. This is why we can easily use parametric equations to solve for the equations of lines that would yield those distances.

 
-WhatdoHumanitariansEat

 \textbf{Alternate finish} 
                    \begin{figure}[H]
                        \centering
                        \includegraphics[width=0.5\textwidth]{Solution/AIME_2000_I_13_Solution_Figure_3.png}
                    \end{figure}
                    Once we get that the area consists of 8 right triangles with side lengths $1.4$, $4.8$, and $5$, then we can finish and find the area in an alternate way which uses no coordinates. Notice how we draw $45^{\circ}$ angles that connect opposite vertices such that the area can be expressed as just the sum of eight congruent triangles. So it suffices to find the area of one of them.

 
                    \begin{figure}[H]
                        \centering
                        \includegraphics[width=0.5\textwidth]{Solution/AIME_2000_I_13_Solution_Figure_4.png}
                    \end{figure}
                    Here we have right triangle $\triangle OAB$ with $\angle OAB=90^{\circ}$, $OA=1.4$, $AB=4.8$, and $OB=5$. $C$ is the point on $AB$ such that $\angle COB=45^{\circ}$. We would like to find the area of $\triangle COB$, so we let $D$ be the foot of the altitude from $C$ to $OB$. We claim the following lemma.

 CLAIM: $AD$ bisects $\angle OAB$

 PROOF: We have $\angle OAC+\angle ODC=90^{\circ}$, so since $\angle OAC+\angle ODC=180^{\circ}$, then $OACD$ is cyclic. Additionally, in $\triangle OCD$, we have $\angle DOC=\angle DCO=45^{\circ}$, so $\angle OAD=\angle CAD=45^{\circ}$. This proves our lemma.

 By the Angle Bisector Theorem, $\frac{OD}{BD}=\frac{AO}{AB}=\frac{7}{24}$. Therefore, $OD=\frac{7}{7+24}OB=\frac{35}{31}$, so $CD=OD=\frac{35}{31}$. Then the area of $\triangle OCB$ is $\frac{1}{2}bh=\frac{1}{2}\cdot 5\cdot \frac{35}{31}=\frac{175}{62}$. The area of the star figure consists of eight of these triangles, so its area is $\frac{175}{62}\cdot 8=\frac{700}{31}$. The answer is $\boxed{731}$.

 




	\item (\textit{AIME I 2000}) In triangle $ABC,$ it is given that angles $B$ and $C$ are congruent. Points $P$ and $Q$ lie on $\overline{AC}$ and $\overline{AB},$ respectively, so that $AP = PQ = QB = BC.$ Angle $ACB$ is $r$ times as large as angle $APQ,$ where $r$ is a positive real number. Find $\lfloor 1000r \rfloor$.


 
 \textbf{Answer: }571



\textbf{Solution 1}


                    \begin{figure}[H]
                        \centering
                        \includegraphics[width=0.5\textwidth]{Solution/AIME_2000_I_14_Solution_Figure_1.png}
                    \end{figure}
                    Let point $R$ be in $\triangle ABC$ such that $QB = BR = RP$. Then $PQBR$ is a \href{https://artofproblemsolving.com/wiki/index.php?title=Rhombus}{rhombus}, so $AB \parallel PR$ and $APRB$ is an \href{https://artofproblemsolving.com/wiki/index.php?title=Isosceles\_trapezoid}{isosceles trapezoid}. Since $\overline{PB}$ bisects $\angle QBR$, it follows by symmetry in trapezoid $APRB$ that $\overline{RA}$ bisects $\angle BAC$. Thus $R$ lies on the perpendicular bisector of $\overline{BC}$, and $BC = BR = RC$. Hence $\triangle BCR$ is an \href{https://artofproblemsolving.com/wiki/index.php?title=Equilateral\_triangle}{equilateral triangle}.

 Now $\angle ABR = \angle BAC = \angle ACR$, and the sum of the angles in $\triangle ABC$ is $\angle ABR + 60^{\circ} + \angle BAC + \angle ACR + 60^{\circ} = 3\angle BAC + 120^{\circ} = 180^{\circ} \Longrightarrow \angle BAC = 20^{\circ}$. Then $\angle APQ = 140^{\circ}$ and $\angle ACB = 80^{\circ}$, so the answer is $\left\lfloor 1000 \cdot \frac{80}{140} \right\rfloor = \left\lfloor \frac{4000}{7} \right\rfloor = \boxed{571}$.

 





\textbf{Solution 2 (Law of sines)}

Let $AP=PQ=QB=BC=x$ and $A$ be the measure of $\angle BAC$. Since $\triangle APQ$ and $\triangle ABC$ are isoceles, $\angle APQ = 180-2A$ and $\angle ACB = 90-\frac{A}{2}$.Because $\triangle APQ$ and $\triangle ABC$ both have a side length $x$ opposite $\angle BAC$, by the law of sines:

 $\frac{x}{\sin A}=\frac{AQ}{\sin(180-2A)}=\frac{AQ+x}{\sin(90-\frac{A}{2})}$

 Simplifying, this becomes

 $\frac{x}{\sin A}=\frac{AQ}{\sin 2A}=\frac{AQ+x}{\cos \frac{A}{2}}$

 From the first two fractions,

 $AQ\cdot \sin A = x \cdot \sin 2A = x \cdot (2\sin A \cos A) \Longrightarrow AQ=x\cdot 2\cos A$

 Substituting, we have from the first and third fractions,

 $\frac{x}{\sin A}=\frac{x\cdot 2\cos A + x}{\cos \frac{A}{2}} \Longrightarrow 2\cos A\sin A + \sin A=\sin 2A + \sin A = \cos \frac{A}{2}$

 By sum-to-product,

 $\sin 2A + \sin A = 2\sin \frac{3A}{2} \cos \frac{A}{2}$

 Thus,

 $2\sin \frac{3A}{2} \cos \frac{A}{2} = \cos \frac{A}{2} \Longrightarrow \sin \frac{3A}{2} = \frac{1}{2}$

 Because $BC=QB<AB$, $\angle A$ is acute, so $\frac{3A}{2}=30 \Longrightarrow A=20$

 $\angle ACB = \frac{180-20}{2}=80$, $\angle APQ = 180-2\cdot 20 = 140 \Longrightarrow r=\frac{4}{7}$

 $1000r=\frac{4000}{7}=\boxed{571}.\overline{428571}$

 


 





\textbf{Solution 3}


                    \begin{figure}[H]
                        \centering
                        \includegraphics[width=0.5\textwidth]{Solution/AIME_2000_I_14_Solution_Figure_2.png}
                    \end{figure}
                    Again, construct $R$ as above.

 Let $\angle BAC = \angle QBR = \angle QPR = 2x$ and $\angle ABC = \angle ACB = y$, which means $x + y = 90$. $\triangle QBC$ is isosceles with $QB = BC$, so $\angle BCQ = 90 - \frac {y}{2}$.Let $S$ be the intersection of $QC$ and $BP$. Since $\angle BCQ = \angle BQC = \angle BRS$, $BCRS$ is \href{https://artofproblemsolving.com/wiki/index.php?title=Cyclic\_quadrilateral}{cyclic}, which means $\angle RBS = \angle RCS = x$.Since $APRB$ is an isosceles trapezoid, $BP = AR$, but since $AR$ bisects $\angle BAC$, $\angle ABR = \angle ACR = 2x$.  

 Therefore we have that $\angle ACB = \angle ACR + \angle RCS + \angle QCB = 2x + x + 90 - \frac {y}{2} = y$.We solve the simultaneous equations $x + y = 90$ and $2x + x + 90 - \frac {y}{2} = y$ to get $x = 10$ and $y = 80$.$\angle APQ = 180 - 4x = 140$, $\angle ACB = 80$, so $r = \frac {80}{140} = \frac {4}{7}$. $\left\lfloor 1000\left(\frac {4}{7}\right)\right\rfloor = \boxed{571}$.

 





\textbf{Solution 4 (Trig identities)}

Let $\angle BAC= 2\theta$ and $AP=PQ=QB=BC=x$. $\triangle APQ$ is isosceles, so $AQ=2x\cos 2\theta =2x(1-2\sin^2\theta)$ and $AB= AQ+x=x\left(3-4\sin^2\theta\right)$. $\triangle{ABC}$ is isosceles too, so $x=BC=2AB\sin\theta$. Using the expression for $AB$, we get 
 \[1=2\left(3\sin\theta-4\sin^3\theta\right)=2\sin3\theta\]by the triple angle formula! Thus $\theta=10^\circ$ and $\angle A = 2\theta=20^\circ$.It follows now that $\angle APQ=140^\circ$, $\angle ACB=80^\circ$, giving $r=\frac{4}{7}$, which implies that $1000r = 571 + \frac 37$. So the answer is $\boxed{571}$.

 





\textbf{Solution 5}


                    \begin{figure}[H]
                        \centering
                        \includegraphics[width=0.5\textwidth]{Solution/AIME_2000_I_14_Solution_Figure_3.png}
                    \end{figure}
                    Reflect $\triangle ABC$ over $BC$ and translate it to attach side $AB$ onto $AC$, mapping $\triangle ABC$ to $\triangle CAD$. Point $P$ maps to $R$, and $Q$ maps to $P$. Then we have that $BC=BQ=QP=PA=AD=PR=RC$. Notice how $BQ=RC$ and $BQ\parallel RC$, so $BQRC$ is a parallelogram and $QR=BC$. But $BC=QP=PR$, so $\triangle QPR$ is actually equilateral. Set $\angle BAC=\angle ACD=x$. Then notice that $\angle QPC=\angle PQA+\angle PAQ=2x$, but $\angle RPC=\angle PQA=x$. Thus $\angle QPR=3x=60$, so $x=20$. Thus $\angle QPA=140^{\circ}$ and $\angle BCA=80^{\circ}$, so $r=\frac{80}{140}=\frac{4}{7}$. The answer is $\left \lfloor \frac{4000}{7}\right \rfloor =\boxed{571}$.

 




	\item (\textit{AIME I 2000}) A stack of $2000$ cards is labelled with the integers from $1$ to $2000,$ with different integers on different cards. The cards in the stack are not in numerical order. The top card is removed from the stack and placed on the table, and the next card is moved to the bottom of the stack. The new top card is removed from the stack and placed on the table, to the right of the card already there, and the next card in the stack is moved to the bottom of the stack. The process - placing the top card to the right of the cards already on the table and moving the next card in the stack to the bottom of the stack - is repeated until all cards are on the table. It is found that, reading from left to right, the labels on the cards are now in ascending order: $1,2,3,\ldots,1999,2000.$ In the original stack of cards, how many cards were above the card labeled $1999$?


 
 \textbf{Answer: }927



\textbf{Solution 1}

We try to work backwards from when there are 2 cards left, since this is when the 1999 card is laid onto the table. When there are 2 cards left, the 1999 card is on the top of the deck. In order for this to occur, it must be 2nd on the deck when there are 4 cards remaining, and this means it must be the 4th card when there are 8 cards remaining. This pattern continues until it is the 512th card on the deck when there are 1024 cards remaining. Since there are over 1000 cards remaining, some cards have not even made one trip through yet, $2(1024 - 1000) = 48$, to be exact. Once these cards go through, 1999 will be the $512 - 48 = 464^\text{th}$ card on the deck. Since every other card was removed during the first round, it goes to show that 1999 was in position $464 \times 2 = 928$, meaning that there were $\boxed{927}$ cards are above the one labeled $1999$.

 





\textbf{Solution 2}

To simplify matters, we want a power of $2$. Hence, we will add $48$ 'fake' cards which we must discard in our actual count. Using similar logic as Solution 1, we find that 1999 has position $1024$ in a $2048$ card stack, where the fake cards towards the front.

 Let the fake cards have positions $1, 3, 5, \cdots, 95$. Then, we know that the real cards filling the gaps between the fake cards must be the cards such that they go to their correct starting positions in the $2000$ card case, where all of them are below $1999$. From this, we know that the cards from positions $1$ to $96$ alternate in fake-real-fake-real, where we have the correct order of cards once the first $96$ have moved and we can start putting real cards on the table. Hence, $1999$ is in position $1024 - 96 = 928$, so $\boxed{927}$ cards are above it. - Spacesam

 





\textbf{Solution 3 (Recursion on number of cards below 1999)}

We work backwards. To reverse the process, we must move the bottom card to the top, and add a new number to the top. Let $d_n$ equal the number of cards below 1999 after $n$ process reversals. Reversing one process, our deck only has $2000$, so we reverse again to obtain $1999, 2000$. So, $d_2 = 1$. When $d_{n-1} > 0$, after a process reversal, the bottom card is moved to the top (it goes above $1999$), so we subtract by one to account for this (the addition of the new card doesn't matter since it goes above $1999$), giving us $d_n = d_{n-1} - 1$. So, $d_3 = 0$. Then, when $d_{n-1} = 0$, after a process reversal, $1999$ moves to the top and one more card is added above $1999$. Since at $d_{n}$, $n$ cards have been added, there must be $n - 2$ cards below $1999$. So, $d_4 = 2$. Then, consider all $d_{n-1} = 0$. Then, $d_n = n-2$ as previously stated. After $n-2$ process reversals, we go back to $d_{2n - 2} = n - 2 - (n - 2) = 0$. Then, $d_{2n-1} = 2n - 3$. Next, we can simply calculate $2 \cdot 4 - 1 = 7$, $2 \cdot 7 - 1 = 13$, and so on (which is a very simple bash, as $2000$ is a small number. If you don't want to do this, define sequence $a_n = 2a_{n-1} - 1$, and solve for the closed form, which is very easy). Consequently, we derive $d_{1537} = 1537 - 2 = 1535$. $2000 - 1537 = 463$ steps remain. So, $d_{2000} = d_{1537} - 463 = 1072$. Finally, there are $2000 - 1 - 1072 = \boxed{\textbf{927}}$ numbers above $1999$.

 





\textbf{Solution 4 (Recursion)}

Consider the general problem: with a stack of $n$ cards such that they will be laid out $1, 2, 3, ..., n$ from left to right, how many cards are above the card labeled $n-1$?

 Let $a_n$ be the answer to the above problem.

 As a base case, consider $n=2$. Clearly, the stack, from top to bottom, must be $(1, 2)$, so $a_2=0$. 

 Next, let's think about how we can construct a stack of $n+1$ cards from a stack of $n$ cards. First, let us renumber the current stack of $n$ by adding $1$ to the label of each of the cards. Then we must add a card labeled "$1$" somewhere in the new deck.

 Working backwards, we find that we must move the bottom card to the top, then add "$1$" to the top of the deck. This is because after one move, the number "$1$" will be laid out and the top card will be moved to the bottom, so the deck becomes the same as what we had before, but with everything renumbered correctly such that $2, 3, 4, ...$ will then be laid out in order.

 Therefore, if $a_n \ne n-1$ (meaning the card $n-1$ is not at the bottom of the deck and so it won't be moved to the top), then $a_{n+1} = a_n + 2$, since a card from the bottom is moved to be above the $n-1$ card, and the new card "$1$" is added to the top. If $a_n = n-1$ (meaning the card $n-1$ is the bottom card), then $a_{n+1}=1$ because the $n-1$ card will move to the top and the card "$1$" will be added on top of it.

 With these recursions and the base case we found earlier, we calculate $a_{2000} = \boxed{927}$. To calculate this by hand, a helpful trick is finding that if $a_n=1$, then $a_{2n-1}=1$ as well. Once we find $a_{1537}=1$, the answer is just $1+(2000-1537)\cdot2$.- Frestho

 





\textbf{Solution 5 (very bashy)}

Let us treat each run through the deck as a separate "round". For example, in round one, you would go through all of the $2000$ cards initially in the deck once, in round two, you would go through all $1000$ cards initially in the deck once, so on and so forth. For each round, let us record what the initial and final actions are ($r$ for moving the card to the right, $b$ for moving the card to the bottom), the number of cards moved to the right, the number of cards left, and what the position of the cards moved to the right were in the original $2000$ card deck (as $n = a + ck$ where $n$ is the position, $c$ is the spacing of the cards moved, $a$ is an integer such that the correct first card is moved, and $k$ is an integer greater than or equal to $1$ which represents which card out of all the cards moved to the right you are finding the position of). Then,

 Round 1: $r$ to $b$, $1000$ to right, $1000$ left in deck, $n = -1 + 2k$,

 Round 2: $r$ to $b$, $500$ to right, $500$ left in deck, $n = -2 + 4k$,

 Round 3: $r$ to $b$, $250$ to right, $250$ left in deck, $n = -4 + 8k$,

 Round 4: $r$ to $b$, $125$ to right, $125$ left in deck, $n = -8 + 16k$.

 Let us treat the remaining deck of $125$ cards as a totally independent deck, note that the positions of card in this deck are $n = 16k$. Also note that the first action of a new round is never the same action as the last action of the previous round because actions alternate. Also note that for every new round, the spacing between cards moved doubles, and that the cards remaining in the beginning of a new round have position $n = a + c/2 + ck$ for the values $a, c$ of the previous round. Also note that if there are an odd number of cards in an initial deck, the first and last actions are the same. Then,

 Round 5: $r$ to $r$, $63$ to right, $62$ left in deck, $n = -1 + 2k$,

 Round 6: $b$ to $r$, $31$ to right, $31$ left in deck, $n = 4k$, because $n = 2k$ positioned cards are left at the beginning of this round and the first card is sent to the bottom, only every second card is sent to the right, and because spacing doubles every round,

 Round 7: $b$ to $b$, $15$ to right, $16$ left in deck, $n = -2 + 8k$, because $n = 2 + 4k$ positioned cards are left at the beginning and only every second card is sent to the right, and because spacing doubles every round,

 Round 8: $r$ to $b$, $8$ to right, $8$ left in deck, $n = -14 + 16k$, by similar reasoning, since the cards $n = 2 + 8k$ are left and spacing doubles every round, from here on things get real easy,

 Round 9: $r$ to $b$, $4$ to right, $4$ left in deck, $n = -22 + 32k$,

 Round 10: $r$ to $b$, $2$ to right, $2$ left in deck, $n = -38 + 64k$,

 Round 11: $r$ to $b$, $1$ to right, $1$ left in deck, $n = 58$, since $58 = -38 + 64/2 + 64$.

 This card must be labeled 1999 since it is second to last. Then, since it is the $58th$ in the deck of $125$, it must be the $58 * 8 = 928th$ card in the original deck, and because we've been numbering from the top of the deck to the bottom (also because AIME answers are 000-999), there were $928 - 1 = \boxed{927}$ cards above it in the deck - Romulus and minimaxweii

 





\textbf{Solution 6 (Bashing with Modular Arithmetic)}

Similar to Solution 5, we treat each run-through of the deck from the lowest-indexed card to the highest-indexed card as a separate round. Notice that after each round, approximately half the deck will remain, with the other half having been cast aside to the right in sorted order. Then, we can model each round as if we are running a "Sieve" through the deck with powers of 2s, filtering out the cards to put to the right and the cards to be pushed to the bottom.

 At the beginning, when all $2000$ cards are still in the deck, notice that the first run-through of the deck consists of removing all the cards such that their index $I\equiv 1\pmod{2}$, while we leave behind all the cards such that their index $I\equiv 2\pmod{2}$. (Since the cards are not $0$-indexed, using $2$ instead of $0\pmod{2}$ is simpler to keep track of the cards. This will appear again later.) The remaining $1000$ cards, then, will all share the attribute that their index numbers are even. To split this further for round 2, we look more specifically at each index, categorizing them further as $I\equiv 2\pmod 4$ or $I\equiv 4\pmod 4$ in that order (The card with the smaller remainder will always be the card on the top of the new deck). Then, in Round 2, we will remove all the cards with index values that leave a remainder of $2$ when divided by $4$, while leaving the multiples of $4$ behind. Notice that each time, in each new deck, we are removing all the cards with an odd position in that specific deck, and leaving behind all the cards with an even position in that specific deck.

 This process will continue until we reach  Round 5, where the number of cards that remain ($125$) is no longer a multiple of 2. Therefore, when we remove all the cards with indices $I\equiv 16\pmod{32}$ and leave behind all the cards with indices $I\equiv 32\pmod{32}$, the last card in the deck will be removed instead of being placed to the bottom, with $125-63=62$ cards remaining. Because of this, when Round 6 begins, we cannot follow our normal procedure, since the top card of the Round 6 deck would have been moved to the bottom! Therefore, when we are using our sieve with $I\equiv 32\pmod{64}$ and $I\equiv 64\pmod{64}$, the order at which the remainders are presented have flipped, so all the cards with indices that are multiples of 64 will be removed, while all the cards with indices that are 32 more than a multiple of 64 will remain. The $62$ cards from Round 6 will be reduced to $31$, and the last card in the deck will be removed just like in Round 5, which means Round 7 will be affected the same way. Of the choices of $I\equiv 32\pmod{128}$ and $I\equiv 32+64=96\pmod{128}$, we will remove the latter, leaving behind the cards with indices that are 32 more than a multiple of 128. However, since 31 is an odd number, the last card of the deck will NOT be removed, which means our sieving process will return to normal after Round 7, with $31-15=16$ cards remaining. Since $16$ is a perfect power of $2$, we will not need to worry about this scenario again in the following rounds.

 Round 8: $I\equiv 32\pmod{256}$ removed; $I\equiv 32+128=160\pmod{256}$ remains; $16-8=8$ cards left.

 Round 9: $I\equiv 160\pmod{512}$ removed; $I\equiv 160+256=416\pmod{512}$ remains; $8-4=4$ cards left.

 Round 10: $I\equiv 416\pmod{1024}$ removed; $I\equiv 416+512=928\pmod{1024}$ remains; $4-2=2$ cards left.

 Round 11: $I\equiv 928\pmod{2048}$ removed; $I\equiv 928+1024=1952\pmod{2048}$ remains; $2-1=1$ card left.

 Round 12: The last card in position $1952$ is sorted into place (This is card \#$2000$!).

 It is obvious that the single card put into place in Round 11 is the second-to-last card in the deck, which is our target \#$1999$. Then, its position in the original deck leaves a remainder of $928$ when divided by $2056$. It is easy to see that $928$ is the only index that satisfies this, so there will be $928-1=\boxed{927}$ cards above card \#$1999$.

 




\end{enumerate}